% ---------------------------------------------------------------------
%  Chapter 4 : Theory Files and Serialization
% ---------------------------------------------------------------------
\chapter{Theory Files and Serialization}\label{ch:theory-files}

The previous chapter introduced the four verbs that drive the engine, and the
first of them --- \code{load\_theory} --- read a small Python file off disk to
get its model. This chapter is about that file: what a \file{.theory.py} is,
exactly what it may contain, how it becomes the in-memory \emph{model dict} the
engine consumes, and the machinery that converts it back and forth between the
graphical theory-builder form and clean, version-controllable source text.

A theory file is the \emph{single source of truth} for one physical model. You
never edit the engine to add a model; you write one of these declarative
modules, and everything downstream --- the field-theory expansion, the
propagator, the diagram enumeration, the loop integrals --- consumes the dict
its \code{build()} produces. Because the file is the authoritative description,
it pays to understand it from the ground up: every chainable builder method,
every module-level constant, and the two distinct ``serialization'' layers that
unfortunately share a name and a directory.

We will work almost entirely from the running example of
Chapter~\ref{ch:quickstart}, \file{theories/ou\_quartic.theory.py}, with a
couple of richer files (a coloured-noise sibling and a spatial KPZ model) brought
in to show the features the simple file does not exercise.

\begin{note}
Two things in this subsystem are both called ``serialization,'' and conflating
them is the single most common source of confusion. The \emph{first} is the
\file{.theory.py} \emph{text} format --- plain, human-readable Python source,
documented in this chapter and handled by \file{pipeline/theory\_serialize.py}.
The \emph{second} is the \file{saved\_theories/<slug>/} \emph{binary} cache ---
pickled SageMath objects (\code{.sobj}) plus a \code{manifest.json}, written when
a theory is actually \emph{run}. You never need Sage to read a \file{.theory.py};
you always need Sage to load a \code{.sobj}. The last section of the chapter
disentangles the two; until then, ``theory file'' always means the text format.
\end{note}

\section{Why theory files exist}

The design promise of a theory file is twofold. First, it is the \term{single
source of truth}: one file fully describes one model --- its fields, parameters,
\msrjd{} action, mean-field equations, and (for spatial models) the boundary and
initial conditions --- with no per-model code anywhere in the engine. Second, it
is the \term{authoring and persistence boundary}: the same content can be
produced either by a human typing the file directly, or by a graphical form in a
Jupyter notebook, and it can be read back into that form for editing.

That second promise is what \file{pipeline/theory\_serialize.py} delivers. The
module is a \emph{bidirectional bridge} between a flat ``spec'' dictionary --- the
internal data the UI form collects --- and \file{.theory.py} source text:

\begin{itemize}
  \item \textbf{Forward} (\code{render\_theory\_file} / \code{save\_theory\_to\_file}):
        the theory-builder UI collects your input as a flat dict and emits a
        clean, diff-friendly \file{.theory.py} file.
  \item \textbf{Reverse} (\code{load\_spec\_from\_file}): given an existing
        \file{.theory.py}, parse its source \emph{as text} --- via Python's
        \code{ast} module, \emph{not} by executing it --- and reconstruct the
        original spec dict, so the UI's ``Load existing theory'' button can
        re-populate the form.
\end{itemize}

The arrangement is worth a diagram. There are two ways a theory file is born (a
UI form, or a human), one canonical way it is loaded (\code{load\_theory}), and a
separate binary cache that is written only when the model is run:

\begin{center}
\begin{tikzpicture}[
    >=Stealth, node distance=6mm,
    box/.style={draw, rounded corners=2pt, align=center,
                font=\footnotesize, inner sep=4pt, fill=codebg},
    file/.style={draw, align=center, font=\footnotesize\ttfamily,
                 inner sep=4pt, fill=notebg},
    lbl/.style={font=\scriptsize\itshape, color=cmt}]
  \node[box] (ui)   {theory-builder UI\\(\code{ipywidgets} form)};
  \node[box, right=28mm of ui] (hand) {hand-authored\\\file{.theory.py}};
  \node[box, below=12mm of ui] (render) {\code{render\_theory\_file(spec)}\\emits source text};
  \node[file, below=24mm of $(ui)!0.5!(hand)$] (tf) {theories/<name>.theory.py};
  \node[box, below=10mm of tf] (model) {model dict\\(\code{mod.build()})};
  \node[box, below=10mm of model] (calc) {FieldTheory / compute\_cumulants\\full\_integrator};
  \node[file, right=14mm of calc] (cache) {saved\_theories/<slug>/\\*.sobj + manifest.json};

  \draw[->] (ui)   -- node[lbl,left]{\code{\_collect()}} (render);
  \draw[->] (render) -- (tf);
  \draw[->] (hand) -- (tf);
  \draw[->] (tf)   -- node[lbl,right]{\code{load\_theory}} (model);
  \draw[->] (model) -- (calc);
  \draw[->] (calc) -- node[lbl,above,align=center]{heavy\\symbolic}(cache);
  % reverse load arrow
  \draw[->, dashed] (tf.west) .. controls +(-3,0) and +(-3,0) ..
       node[lbl,left,align=center]{\code{load\_spec}\\\code{\_from\_file}} (render.west);
\end{tikzpicture}
\end{center}

The dashed arrow is the reverse path: it reads the file back into a spec so the
form can show it. The solid path down the middle is the run-time path you met in
Chapter~\ref{ch:quickstart}.

The key correctness property of the forward/reverse pair is \term{round-trip
fidelity}: a cycle \code{spec} $\to$ \code{render} $\to$ \code{file} $\to$
\code{load} $\to$ \code{spec'} reproduces the \emph{meaningful} content of the
original spec. Because the reverse path reads the abstract syntax tree rather
than running the module, string literals (action text, equation right-hand
sides, kernel expressions) round-trip \emph{verbatim}, while comments and exact
formatting are lost. We will return to the precise limits of that guarantee in
\S\ref{sec:tf-reverse}.

\section{Anatomy of a \texttt{.theory.py}}

A theory file is an ordinary, importable Python module with exactly three public
symbols:

\begin{description}
  \item[\code{build()}] a function returning the \term{model dict} --- the
        in-memory object the rest of the pipeline consumes. Internally it
        constructs the model with a fluent \emph{builder} object whose chained
        method calls declare the fields, parameters, action, and equations.
  \item[\code{DEFAULT\_FUNDAMENTAL}] a plain dict of default numeric parameter
        values --- the ``working point'' the calculation is evaluated at.
  \item[\code{METADATA}] a plain dict of \emph{run recommendations}: the default
        cumulant order $k$, the default loop order $\ell$, which external fields
        to plot, the time grid, and (for spatial models) the spatial grid.
\end{description}

Here is the entire white-noise quartic file again, this time as the object of
study rather than a passing reference (\file{theories/ou\_quartic.theory.py},
docstring elided):

\begin{lstlisting}
from pipeline.theory import TemporalTheoryBuilder


def build():
    return (
        TemporalTheoryBuilder('OU Quartic (white noise)')
        .population('pop', size=1)
        .physical_field('x', population='pop', description='variable')
        .parameter('mu', default=1.0, domain='positive')
        .parameter('eps', default=0.02, domain='positive')
        .parameter('D', default=1.0, domain='positive')
        .set_action_text(
            'sum(xt[i]*((Dt+mu)*x[i] + eps*x[i]^3) - D*xt[i]^2 for i in pop)')
        .equation(lhs='(Dt+mu)*x[i]', rhs='-eps*x[i]^3', population='pop')
        .build()
    )


DEFAULT_FUNDAMENTAL = {'mu': 1.0, 'eps': 0.02, 'D': 1.0}

METADATA = {
    'k_default': 2,
    'ell_default': 0,
    'recommended_external_fields': [('dx', 1), ('dx', 1)],
    'tau_max': 8.0,
    'tau_step': 0.5,
    'fixed_point_index_default': 0,
}
\end{lstlisting}

Read top to bottom this is almost self-documenting. The \code{import} line names
the builder class (here \code{TemporalTheoryBuilder}; \S\ref{sec:tf-split}
explains the temporal/spatial choice). \code{build()} returns a single
parenthesised expression: a constructor call followed by a chain of
\code{.method(...)} calls, terminated by \code{.build()}. Each method declares
one ingredient of the model. \code{DEFAULT\_FUNDAMENTAL} and \code{METADATA} are
module-level constants the loader hands back alongside the model dict --- they are
\emph{not} part of \code{build()}'s return value, which is why
\code{load\_theory} returns the \emph{pair} \code{(model, module)} you saw in
Chapter~\ref{ch:quickstart}: the model dict carries the physics, the module
carries the run recommendations.

\subsection{The math the file encodes}

A theory file encodes a stochastic field theory in the \msrjd{} formalism. We
assume you know that physics; here is only the minimal vocabulary needed to read
the code.

Each physical field \code{x} (or \code{phi}, \code{h}, \code{n}, \dots) gets a
conjugate \term{response field} \code{xt} (the ``tilde'' field $\tilde x$). A
Langevin process
\[
  \frac{\dd x}{\dd t} = F[x] + \xi,
  \qquad \avg{\xi(t)\,\xi(t')} = 2D\,\delta(t-t'),
\]
is rewritten as a path integral with the \msrjd{} action $S[x,\tilde x]$. For the
running example,
\[
  \frac{\dd x}{\dd t} = -\mu x - \varepsilon x^{3} + \xi,
\]
the action --- exactly as it appears in \code{set\_action\_text} --- is
\begin{equation}
  S \;=\; \int \dd t \,\Bigl[\, \tilde x\bigl((\Dt + \mu)\,x + \varepsilon x^{3}\bigr)
       \;-\; D\,\tilde x^{2} \,\Bigr].
  \label{eq:tf-action}
\end{equation}
The string passed to \code{.set\_action\_text(...)} is this expression with
\code{xt} for $\tilde x$ and \code{\^{}} for powers
(\file{theories/ou\_quartic.theory.py:27--28}):

\begin{lstlisting}
'sum(xt[i]*((Dt+mu)*x[i] + eps*x[i]^3) - D*xt[i]^2 for i in pop)'
\end{lstlisting}

The time-derivative $\partial_t$ is the operator \code{Dt}; the quadratic
response term $-D\,\tilde x^{2}$ is the white-noise vertex; the cubic
$\varepsilon x^{3}$ is the single interaction. The \code{sum(\dots\ for i in pop)}
wrapper is a \emph{population comprehension} --- see \S\ref{sec:tf-indexing}.

\paragraph{Mean field / saddle point.} The classical (noise-free) trajectory is
the \term{mean-field saddle} $x^{*}$. Theory files declare it one of two ways.
The preferred, newer form is the \term{DAE} (differential-algebraic equation)
\code{.equation(lhs=..., rhs=..., population=...)}; in the running example
(\file{theories/ou\_quartic.theory.py:29})

\begin{lstlisting}
.equation(lhs='(Dt+mu)*x[i]', rhs='-eps*x[i]^3', population='pop')
\end{lstlisting}

reads as ``$(\partial_t + \mu)\,x = -\varepsilon x^{3}$,'' a DAE the solver
integrates to a fixed point. The legacy alternative names the saddle symbol
explicitly with \code{.set\_mf\_equation('xstar', 'rhs')}; you will meet it in
heterogeneous Hawkes models. The forward serializer prefers \code{.equation(...)}
when the spec carries an \code{equations} list and falls back to
\code{.set\_mf\_equation(...)} otherwise
(\file{pipeline/theory\_serialize.py:424--433}).

\paragraph{Propagator, cumulants, loop expansion.} After choosing $x^{*}$, the
theory is expanded around it. The \term{propagator} is the inverse of the
quadratic part of $S$; \term{cumulants} (connected correlation functions) are
computed order by order in a \term{loop expansion}. Two integers parametrise the
request, and you met both as \code{Config} fields in Chapter~\ref{ch:quickstart}:
$k$, the cumulant order (number of external legs), and $\ell$ (the
\code{max\_ell} / loop order), the highest loop order included. In the theory
file these appear in \code{METADATA} as \code{k\_default} and \code{ell\_default},
and they later become file-name keys in the binary cache (e.g.
\code{...\_k2\_l1.sobj}).

\begin{note}
The small parameter of the loop expansion for this model is roughly
$g_{\mathrm{eff}} \approx \varepsilon D / \mu^{2}$. The comment at
\file{theories/ou\_quartic.theory.py:34} records $g_{\mathrm{eff}} \approx 0.02$
at the default working point, which is why the series converges fast. This is the
same number that drove the regime discussion in Chapter~\ref{ch:quickstart}; the
theory file is simply where that working point is written down.
\end{note}

\subsection{\texttt{DEFAULT\_FUNDAMENTAL} and \texttt{METADATA}}

\code{DEFAULT\_FUNDAMENTAL} is the numeric \emph{working point}. It may contain
scalars, vectors (lists), or matrices (lists of lists). Three styles appear
across the shipped theories:

\begin{itemize}
  \item \code{\{'mu': 1.0, 'eps': 0.02, 'D': 1.0\}} --- scalar defaults
        (\code{ou\_quartic}).
  \item \code{\{\}} --- \emph{empty}, when defaults are carried on each
        \code{.parameter(default=...)} instead (as in \code{ou\_quartic\_colored}
        and \code{quadratic\_hawkes\_alpha}).
  \item \code{\{'E': [0.78, 0.81], 'w': [[0.30,0.25],[0.30,0.35]], ...\}} ---
        vector and matrix entries for heterogeneous populations
        (\code{linear\_hawkes}).
\end{itemize}

\begin{gotcha}
\code{DEFAULT\_FUNDAMENTAL} \emph{can be empty}. Several theories leave it
\code{\{\}} and carry the numbers on the individual \code{.parameter(default=...)}
calls. Both styles are valid; \code{run} merges the model's own parameter
defaults, then \code{DEFAULT\_FUNDAMENTAL}, then your \code{Config.parameters}
(the layering you saw in Chapter~\ref{ch:quickstart}), so a freshly loaded theory
runs out of the box regardless of which style its author chose.
\end{gotcha}

\code{METADATA} is the bundle of \emph{run recommendations}. The keys observed
across the shipped theories are:

\begin{center}
\begin{longtable}{@{}llp{0.46\textwidth}@{}}
\toprule
\textbf{Key} & \textbf{Example} & \textbf{Meaning} \\
\midrule
\endhead
\code{k\_default} & \code{2} & Default cumulant order. \\
\code{ell\_default} & \code{0}, \code{1} & Default loop order. \\
\code{recommended\_external\_fields} & \code{[('dx',1),('dx',1)]} &
  \code{(field, population\_index)} legs to compute/plot. \\
\code{tau\_max}, \code{tau\_step} & \code{8.0}, \code{0.5} & Time grid. \\
\code{fixed\_point\_index\_default} & \code{0} & Which saddle root to use
  (temporal only). \\
\code{spatial\_grid} & \code{[-8.0, 8.0, 65]} & Spatial sampling, as an explicit
  list or a \code{(start, stop, n)} triple. \\
\code{description} & string & Optional short label. \\
\bottomrule
\end{longtable}
\end{center}

\begin{gotcha}
The entries of \code{recommended\_external\_fields} are \emph{tuples}, e.g.
\code{('dx', 1)}, and downstream code expects the shape
\code{(field\_name, population\_index)}. They render as tuples and reload as
tuples (the reverse path recovers them with \code{ast.literal\_eval}). If you
hand-author \code{METADATA} and write them as \emph{lists}, the type would change
on reload --- so keep them tuples.
\end{gotcha}

\section{The builder DSL}

The body of \code{build()} is a fluent chain on a builder object. Each method is
a small declaration; the chain reads top to bottom like a specification. This
section walks every chainable method you are likely to write, grouped by what it
declares. The methods live on \code{TemporalTheoryBuilder} /
\code{SpatialTheoryBuilder} (imported from \code{pipeline.theory}); the forward
serializer's \code{\_emit\_*} helpers, cited per method below, are the authority
on exactly which keyword arguments each call accepts.

\subsection{Fields and parameters}

\paragraph{\code{.physical\_field('x', \dots)}.} Declares one physical field. The
positional \code{'x'} is the user-facing \emph{natural name}; the builder
auto-prefixes \code{d} to form the internal fluctuation name (\code{dx}) and
auto-creates both the conjugate response field (\code{xt}) and the saddle
(\code{xstar}). This is why the model in Chapter~\ref{ch:quickstart} reported its
field as \code{dx} even though the file wrote \code{x}, and why you almost never
write a \code{.response\_field(...)} line by hand. Optional keywords, emitted by
\code{\_emit\_field} (\file{pipeline/theory\_serialize.py:112}), are
\code{population=} (heterogeneous style) \emph{or} \code{indexed=} (legacy
style), \code{natural\_name=}, \code{latex=}, \code{description=}, and a per-field
\code{spatial\_dim=}.

\paragraph{\code{.parameter('mu', default=1.0, domain='positive')}.} Declares one
numeric parameter. The keywords (emitted by \code{\_emit\_parameter},
\file{pipeline/theory\_serialize.py:151}) are \code{default=}, the indexing
declaration (\code{indexed\_by=} new-style or \code{indexed=} legacy),
\code{domain=} (e.g. \code{'positive'}, \code{'real'}), \code{natural\_name=}, and
\code{description=}.

\begin{gotcha}
\textbf{Saddle parameters are never written by hand and never emitted.} A
parameter whose name matches the auto-generated \code{<natural>star} convention
(e.g. \code{xstar}), or one flagged \code{mean\_field}, is \emph{filtered out} by
the forward serializer (\file{pipeline/theory\_serialize.py:334--339}) because the
framework re-creates it from the physical-field declaration. Emitting it would
double-declare the saddle. The filter is even applied twice --- belt and braces ---
since \code{\_emit\_parameter} guards \code{mean\_field} again at line~181.
\end{gotcha}

\paragraph{\code{.response\_field('xt', \dots)}.} Rarely needed --- only emitted
when you \emph{explicitly} declared a response field rather than relying on the
auto-creation a \code{.physical\_field(...)} performs. Most files therefore carry
no \code{.response\_field(...)} line at all.

\subsection{Functions and kernels}

\paragraph{\code{.define\_function('phi', args=[...], expression='...')}.} Declares
a named scalar function used in the action or saddle (handler
\code{\_emit\_function}, \file{pipeline/theory\_serialize.py:189}). \code{args}
and \code{expression} are mandatory; \code{population=}, \code{latex=}, and
\code{description=} are optional. For example
\code{.define\_function('phi', args=['v'], expression='a[i]*v\^{}2', population='E')}
declares a per-population quadratic transfer function. Functions exist precisely
to support the saddle indirection some models need.

\paragraph{\code{.define\_kernel(\dots)}.} Declares a convolution kernel (handler
\code{\_emit\_kernel}, \file{pipeline/theory\_serialize.py:200}). A kernel may be
given in the time domain via \code{time\_expr=} (e.g.
\code{'(t/taug[i,j]\^{}2)*exp(-t/taug[i,j])*heaviside(t)'}) \emph{or} directly as
a frequency image via \code{freq\_image=} (e.g.
\code{'1 / (1 + I*omega*tau\_g)'}). It also accepts \code{latex\_name=},
\code{sage\_name=}, \code{description=}, and the same indexing duality as
parameters.

\subsection{The action and the mean-field equations}

\paragraph{\code{.set\_action\_text('...')}.} The heart of the model: the \msrjd{}
action of \eqref{eq:tf-action} as a single string, with \code{xt} for the
response field and \code{\^{}} for powers. The handler \code{\_emit\_action}
(\file{pipeline/theory\_serialize.py:277}) emits single-line bodies as a plain
\code{repr}; multi-line bodies are first run through \code{textwrap.dedent} and
then re-indented four spaces inside a triple-quoted block. The dedent is
load-bearing for whitespace stability --- see the gotcha in \S\ref{sec:tf-forward}.

\paragraph{\code{.equation(lhs=..., rhs=..., population=...)} (preferred) and
\code{.set\_mf\_equation('xstar', 'rhs')} (legacy).} The two saddle-declaration
forms introduced above. The new DAE form is emitted by \code{\_emit\_equation}
(\file{pipeline/theory\_serialize.py:297}), which omits \code{population=} when it
is falsy; the legacy form by \code{\_emit\_mf\_equation}
(\file{pipeline/theory\_serialize.py:292}).

\subsection{Toggles}

A handful of methods set boolean toggles or scalar policies. The forward
serializer emits each \emph{only when it departs from the builder's default}, so
files stay ``textually quiet'' --- a recurring design phrase meaning a file
contains only what is non-trivial, which keeps diffs clean. These are:

\begin{itemize}
  \item \code{.stability\_analysis(True)} --- emitted only when on
        (\file{pipeline/theory\_serialize.py:439--440});
  \item \code{.markovianize(False)} --- emitted only when explicitly off
        (\file{pipeline/theory\_serialize.py:448--449}); see
        \S\ref{sec:tf-colored};
  \item \code{.operator\_ir()} --- emitted only when on, and \emph{load-bearing}
        for derivative-vertex spatial theories
        (\file{pipeline/theory\_serialize.py:452--459}); see
        \S\ref{sec:tf-spatial};
  \item \code{.boundary(mode, **params)} / \code{.initial(mode, **params)} ---
        spatial boundary and initial conditions, emitted only when present
        (\file{pipeline/theory\_serialize.py:466--477});
  \item \code{.dyson\_order(N)} / \code{.reference\_diffusion(D0)} --- the
        Dyson--Duhamel policy and reference-diffusion split, emitted only for a
        non-\code{'off'} policy or a set $D_{0}$
        (\file{pipeline/theory\_serialize.py:484--489}).
\end{itemize}

\section{Temporal versus spatial: the builder split}\label{sec:tf-split}

The \code{import} line at the top of a theory file names one of two builder
classes. \code{TemporalTheoryBuilder} is for time-only (ODE) models like the
running example; \code{SpatialTheoryBuilder} is for spatial (PDE) models where
the field lives on $x \in \mathbb{R}^{d}$ and the action carries the Laplacian.

The forward serializer does \emph{not} ask you which one you want --- it infers it.
The choice is driven purely by whether any physical field carries a non-zero
\code{spatial\_dim} (\file{pipeline/theory\_serialize.py:370--372}):

\begin{lstlisting}
_is_spatial = any(int(f.get('spatial_dim', 0) or 0) > 0
                  for f in physical_fields)
builder_cls = 'SpatialTheoryBuilder' if _is_spatial else 'TemporalTheoryBuilder'
\end{lstlisting}

If any field declares $\code{spatial\_dim} \ge 1$, the file is emitted as spatial;
otherwise it is temporal.

\begin{gotcha}
The builder split is \emph{inferred from \code{spatial\_dim}, not declared}. A
spatial theory whose fields all forgot \code{spatial\_dim} would be emitted as a
\emph{temporal} file --- and would then compute the wrong physics. The reverse
(loading) path is forgiving: it accepts \code{TheoryBuilder} (a legacy
back-compat shim), \code{SpatialTheoryBuilder}, \emph{and}
\code{TemporalTheoryBuilder} as constructor names
(\file{pipeline/theory\_serialize.py:651--652}). But the forward path can only
emit the right class if at least one field carries the dimension.
\end{gotcha}

\section{Indexing and populations}\label{sec:tf-indexing}

A theory can be \term{homogeneous} (one population, scalar fields) or
\term{heterogeneous} (multiple populations of different sizes, with vector or
matrix parameters). The action's \code{sum(... for i in pop)} comprehension you
saw above is how a field equation is written once and applied across a
population: \code{i} ranges over the population \code{pop}, and \code{x[i]},
\code{xt[i]} are the per-member fields. For the running example \code{pop} has
size~1, so the sum is a single term --- but the same syntax scales to genuine
networks.

Indexing is encoded in one of two styles, and the serializer round-trips both:

\begin{description}
  \item[New population style.] \code{.population('E', size=2)} declares a named
        population of size $N$; fields attach via \code{population='E'} and
        parameters via \code{indexed\_by=['E']} (vector) or
        \code{indexed\_by=['E','E']} (matrix). The running example uses the
        single-member form \code{.population('pop', size=1)}.
  \item[Legacy style.] \code{n\_populations=N} in the constructor, plus
        \code{indexed=True} (vector) or \code{indexed='matrix'} on parameters. The
        coloured-noise file in \S\ref{sec:tf-colored} uses
        \code{TemporalTheoryBuilder('OU Quartic Colored', n\_populations=0)} with
        \code{.physical\_field('x', indexed=True)}.
\end{description}

\begin{gotcha}
\code{indexed\_by} \emph{wins} over the legacy \code{indexed}/\code{type}. Both
\code{\_emit\_parameter} and \code{\_emit\_kernel} accept either; if a spec
somehow carries both, the new population-style \code{indexed\_by} is used and the
legacy \code{indexed=} is suppressed. Mixing the two styles in a single
hand-written file is best avoided --- pick one.
\end{gotcha}

\section{Colored noise and CGF terms}\label{sec:tf-colored}

White noise enters the action as the local term $-D\,\tilde x^{2}$ --- a single
noise vertex with no memory. \term{Colored} (finite-correlation-time) noise
instead enters as a \term{CGF term} (cumulant-generating-function term), declared
with \code{.declare\_cgf\_term(...)} and carrying an \code{order}, a
\code{coefficient}, and a temporal \code{kernel}. The coloured sibling of the
running model, \file{theories/ou\_quartic\_colored.theory.py}, is the canonical
example. Its \code{build()} reads (verbatim, \file{...:10--23}):

\begin{lstlisting}
def build():
    return (
        TemporalTheoryBuilder('OU Quartic Colored', n_populations=0)
        .physical_field('x', indexed=True, description='variable')
        .parameter('mu', default=0.1, domain='real')
        .parameter('eps', default=0.1, domain='positive')
        .parameter('D', default=1.0, domain='positive')
        .parameter('tauc', default=1.0, domain='positive')
        .declare_cgf_term('CXX', response_legs=['xt', 'xt'], order=2,
                          coefficient='2*D/tauc', kernel='exp(-abs(tau)/tauc)')
        .set_action_text('xt*((Dt+mu)*x + eps*x^3)')
        .equation(lhs='(Dt+mu)*x', rhs='-eps*x^3')
        .stability_analysis(True)
        .build()
    )
\end{lstlisting}

Two differences from the white-noise file are the whole point. First, the action
\code{'xt*((Dt+mu)*x + eps*x\^{}3)'} has \emph{no} $-D\,\tilde x^{2}$ term --- the
noise is no longer local. Second, that noise is supplied by the CGF row

\begin{lstlisting}
.declare_cgf_term('CXX', response_legs=['xt', 'xt'], order=2,
                  coefficient='2*D/tauc', kernel='exp(-abs(tau)/tauc)')
\end{lstlisting}

which encodes the Ornstein--Uhlenbeck noise spectrum
$\avg{\xi\,\xi} \propto (D/\tau_c)\,\ee^{-|\tau|/\tau_c}$. The \code{response\_legs}
list says the term couples two response fields ($\tilde x$ with $\tilde x$), the
\code{order} is the cumulant order of the term, \code{coefficient} its prefactor,
and \code{kernel} its time profile.

The handler \code{\_emit\_cgf\_term} (\file{pipeline/theory\_serialize.py:233})
renders two shapes. The multi-leg shape above is used when \code{response\_legs}
has more than one entry; a legacy single-leg shape
\code{.declare\_cgf\_term('name', 'mt', order=..., ...)} passes the single
response field as the second positional argument. \code{order} is forced to
\code{int}, and a per-row \code{markovianize=} is emitted only when it is an
explicit boolean.

\paragraph{Markovianization.} By default the pipeline \term{Markovianizes} a
coloured kernel --- it embeds the finite-correlation-time noise in an auxiliary
Markovian field so the system becomes local in time again (the embedding lives in
\file{pipeline/colored\_to\_markovian.py}). The builder-level
\code{.markovianize(False)} opts the whole theory out, and a per-term
\code{markovianize=True/False} overrides individual rows. The forward serializer
emits the builder-level \code{.markovianize(False)} only when the spec explicitly
turned it off (\file{pipeline/theory\_serialize.py:448--449}), so theories that
accept the default stay quiet.

\begin{note}
Chapter~\ref{ch:quickstart}'s note on ``what colour costs'' is exactly this
mechanism: the white-noise model has no $\tau_c$ to embed, hence no auxiliary
field and a small diagram count, while the coloured sibling adds the embedding.
The two files differ in precisely the action term and the one CGF row shown
above. The dedicated chapter on coloured and Markovian noise
(Chapter~\ref{ch:colored-markovian}) opens up the embedding itself.
\end{note}

\section{Spatial derivative vertices and the operator IR}\label{sec:tf-spatial}

For spatial theories the action carries the Laplacian $\nabla^{2}$ and, for
derivative-vertex models, gradient operators. There are two distinct authoring
styles, and the difference between them is not cosmetic --- it changes the physics
the loop integrals compute.

\begin{description}
  \item[Bare-multiplicative Laplacian (v1).] The inert symbol \code{Laplacian} is
        used multiplicatively, e.g. \code{(Dt + mu - D*Laplacian)*phi}. No
        operator IR is needed.
  \item[Operator IR (v2).] The action uses \emph{call syntax} --- \code{Dt(h)},
        \code{Lap(h)}, \code{Dx(h, 0)} --- for genuine derivative \emph{vertices}.
        A theory using call syntax \emph{must} carry \code{.operator\_ir()}.
\end{description}

\begin{defn}
The \term{operator IR} is the v2 intermediate representation in which derivative
operators are call-syntax nodes (\code{Dt(h)}, \code{Lap(h)}, \code{Dx(h, axis)})
rather than inert multiplicative symbols. Because a derivative \emph{vertex}
contributes a momentum-dependent \term{form factor} to a loop integral, the IR is
what lets the engine attach the right form factor to each vertex. The integer in
\code{Dx(h, 0)} is the \term{spatial axis index} (here $d=1$, so axis~0).
\end{defn}

Why the IR is load-bearing is best seen on KPZ. The Kardar--Parisi--Zhang
equation in \file{theories/kpz\_1d.theory.py} is
\[
  \partial_t h = -\mu h + D\nabla^{2}h + \tfrac{\lambda}{2}(\partial_x h)^{2} + \eta,
  \qquad \avg{\eta\,\eta} = 2T\,\delta(x-x')\,\delta(t-t'),
\]
written as

\begin{lstlisting}
.set_action_text(
    'ht*(Dt(h) + mu*h - D*Lap(h) - (lam/2)*Dx(h, 0)^2) - T*ht^2')
.operator_ir()
.boundary('infinite')
.initial('stationary')
\end{lstlisting}

The KPZ nonlinearity $\tfrac{\lambda}{2}(\partial_x h)^{2}$ is a \emph{per-leg}
gradient vertex, written \code{Dx(h, 0)\^{}2}. Contrast Model B's conserving flux
$g\nabla^{2}(\phi^{2})$, a \emph{composite} derivative vertex written
\code{g*Lap(phi\^{}2)}. These differ --- $(\partial_x h)^{2} \neq
\partial_x(h^{2})$ --- and they produce different momentum form factors in the
loop integral.

\begin{gotcha}
\code{.operator\_ir()} is \textbf{load-bearing and silent}. The comment at
\file{pipeline/theory\_serialize.py:452--457} is explicit: a re-saved
KPZ/Burgers/Model-B theory that loses the \code{.operator\_ir()} line will
``silently lose the \code{Dt()/Lap()/Dx()} lowering and compute the wrong
physics.'' The serializer therefore round-trips the toggle precisely, and emits it
only when on (so v1 bare-Laplacian theories stay quiet). \emph{If you hand-edit a
derivative-vertex theory, keep that line.}
\end{gotcha}

For spatial models you will also see \code{.boundary(mode, **params)} (e.g.
\code{.boundary('infinite')} or a periodic box with a \code{length=}) and
\code{.initial(mode, **params)} (e.g. \code{.initial('stationary')}), and for
coupled multi-field models the \code{.dyson\_order(N)} /
\code{.reference\_diffusion(D0)} policy pair. The spatial chapters of Part~V cover
what those declarations route to.

\section{The spec dict: the contract with the UI}\label{sec:tf-spec}

Everything above describes the \emph{file}. Internally, the UI form and the
serializer speak a flat dictionary --- the \term{spec dict}. It is produced by the
form's \code{\_collect()}, rendered to text by \code{render\_theory\_file}, and
reconstructed from text by \code{load\_spec\_from\_file}. Knowing its shape makes
both the forward and reverse paths legible, because every \code{\_emit\_*} helper
is really ``read this key, write that line,'' and every reverse-path branch is the
mirror image. The defaults below are exactly those the reverse path initialises
(\file{pipeline/theory\_serialize.py:601--637}):

\begin{center}
\begin{longtable}{@{}llp{0.50\textwidth}@{}}
\toprule
\textbf{Key} & \textbf{Type} & \textbf{Meaning} \\
\midrule
\endhead
\code{name} & str & Human theory name (also the slug source). \\
\code{description} & str & User prose (boilerplate-stripped on reload). \\
\code{populations} & list[dict] & \code{[\{'name','size','description'?\}, ...]}
  (heterogeneous style). \\
\code{n\_populations} & int & Legacy population count; synced to
  \code{len(populations)} when populations are present. \\
\code{response\_fields} & list[dict] & Explicitly-declared response fields;
  usually empty. \\
\code{physical\_fields} & list[dict] & Field records (name, indexing, natural
  name, latex, description, \code{spatial\_dim}). \\
\code{parameters} & list[dict] & Parameter records (name, default, indexing,
  domain, \dots). \\
\code{functions} & list[dict] & \code{define\_function} records (name, args,
  expression, \dots). \\
\code{kernels} & list[dict] & \code{define\_kernel} records (time/freq exprs,
  names, indexing). \\
\code{cgf\_terms} & list[dict] & Coloured-noise rows (legs, order, coefficient,
  kernel, markovianize). \\
\code{action\_text} & str & The \msrjd{} action expression (verbatim). \\
\code{mf\_equations} & list[dict] & Legacy \code{[\{'saddle','rhs'\}, ...]}. \\
\code{equations} & list[dict] & DAE form \code{[\{'lhs','rhs','population'\}, ...]}. \\
\code{stability\_analysis} & bool & Default \code{False}. \\
\code{markovianize\_default} & bool & Default \code{True}. \\
\code{operator\_ir} & bool & Default \code{False}. \textbf{Load-bearing} for
  derivative-vertex spatial theories. \\
\code{boundary} & dict & \code{\{'mode', **params\}}, e.g.
  \code{\{'mode':'infinite'\}}. \\
\code{initial} & dict & \code{\{'mode', **params\}}, e.g.
  \code{\{'mode':'stationary'\}}. \\
\code{dyson} & dict & \code{\{'mode':'fixed','order':N\}}. \\
\code{reference\_diffusion} & float & Scalar $D_{0}$ for the
  $\mathcal{D} = D_{0}\,I + \hat{\mathcal{D}}$ split. \\
\code{default\_fundamental} & dict & Emitted as \code{DEFAULT\_FUNDAMENTAL}. \\
\code{metadata} & dict & Emitted as \code{METADATA}. \\
\bottomrule
\end{longtable}
\end{center}

One transient key, \code{\_default\_spatial\_dim}, is set during reverse parsing
by a \code{.spatial\_dim(d)} call and is not normally re-emitted; it exists only
so a bulk \code{.spatial\_dim(d)} can back-fill fields that lacked a per-field
\code{spatial\_dim=}. We return to it in \S\ref{sec:tf-reverse}.

\section{Forward serialization: spec to file}\label{sec:tf-forward}

The forward heart is \code{render\_theory\_file(spec) -> str}
(\file{pipeline/theory\_serialize.py:310}), which takes a spec dict and returns
the full \file{.theory.py} source string. \code{save\_theory\_to\_file(spec, path)}
(\file{...:504}) is the thin wrapper that picks a filename and writes it. The
module itself is deliberately \emph{dependency-light}: it touches only the Python
standard library (\code{ast}, \code{re}, \code{textwrap}, \code{os}, \code{typing}).
No SageMath, no \code{ipywidgets}, no \code{numpy} --- those enter only in the
companion loader and the binary cache.

\code{render\_theory\_file} proceeds in a fixed order. Abbreviated, it:

\begin{enumerate}
  \item reads the spec fields with \code{.get(key, default)} so partial specs
        still render (\file{...:320--351});
  \item filters out saddle parameters --- those flagged \code{mean\_field} or named
        \code{<natural>star} (\file{...:334--339});
  \item picks the builder class from \code{spatial\_dim} (\S\ref{sec:tf-split},
        \file{...:370--372});
  \item emits the docstring header --- a \code{"<name> --- text-driven theory
        file."} line plus the optional description and the ``Generated by\dots\ /
        Loaded by\dots'' boilerplate the reverse path later strips
        (\file{...:354--363});
  \item emits \code{from pipeline.theory import <BuilderCls>} and
        \code{def build(): return (} (\file{...:376--380});
  \item emits the constructor: with \code{populations} present, a bare
        \code{BuilderCls('name')} followed by one \code{.population(...)} per
        population; otherwise \code{BuilderCls('name', n\_populations=N)}
        (\file{...:381--396});
  \item emits, in fixed order, response fields, physical fields, parameters,
        functions, kernels, CGF terms (\file{...:401--412});
  \item emits the action, then the equations (preferring \code{.equation(...)},
        else \code{.set\_mf\_equation(...)}) (\file{...:414--433});
  \item emits the non-default toggles --- stability, \code{markovianize(False)},
        \code{operator\_ir()}, boundary/initial, Dyson/reference-diffusion
        (\file{...:439--489});
  \item closes the chain with \code{.build()} and the \code{)}
        (\file{...:491--492}), then appends \code{DEFAULT\_FUNDAMENTAL = \{...\}}
        and \code{METADATA = \{...\}} (\file{...:495--498}).
\end{enumerate}

The low-level work is done by a handful of helpers worth knowing by name.
\code{\_py\_repr(value)} (\file{...:46}) pretty-prints a Python value as a
source-code literal: a 2D list (a matrix) becomes a multi-line block with one row
per line, a short 1D list stays on one line, a dict becomes one
\code{key: value} per line. This is what makes a matrix default print as a
readable block rather than one long line. \code{\_kw\_chain(*pairs)}
(\file{...:85}) joins \code{name=value} pairs but \emph{omits any pair whose value
is \code{None} or \code{''}} --- this is the mechanism behind the ``textually
quiet'' principle, so \code{.parameter('mu', default=1.0)} does not emit
\code{description=None}. \code{\_slugify(name)} (\file{...:96}) is the
filesystem-safe identifier
\code{re.sub(r'[\^{}A-Za-z0-9]+', '\_', name).strip('\_').lower()}.

\begin{gotcha}
In \code{\_py\_repr}, \code{bool} is checked \emph{before} \code{int}
(\file{...:54}). In Python \code{bool} is a subclass of \code{int}, so a naive
\code{isinstance(value, int)} test would catch \code{True}/\code{False} first and
print them as \code{1}/\code{0}. The order is deliberate; do not swap it.
\end{gotcha}

\begin{gotcha}
Whitespace stability depends on the dedent in \code{\_emit\_action}
(\file{...:281--289}). The action body is run through \code{textwrap.dedent}
\emph{before} being re-indented four spaces. Without that dedent, every
load$\to$save cycle would accumulate four more leading spaces on each line of a
multi-line action. Do not ``simplify'' it away.
\end{gotcha}

\code{save\_theory\_to\_file} normalises the path before writing
(\file{...:513--518}): if \code{path} is a directory, the filename is
\code{<slugify(name)>.theory.py} inside it; otherwise, if \code{path} does not
already end in \code{.theory.py}, the suffix is appended \emph{unless} it already
ends in \code{.py}.

\begin{gotcha}
The extension logic has a subtlety. A path ending in \code{.py} but \emph{not}
\code{.theory.py} is left untouched, so \code{save\_theory\_to\_file(spec, 'foo.py')}
writes \code{foo.py} --- which \code{dd.list\_theories()} would then \emph{not}
discover, because discovery requires the \code{.theory.py} suffix
(\file{notebooks/daedalus.py:63--65}). If you want a loadable theory, give it a
bare name or the full \code{.theory.py} suffix.
\end{gotcha}

\section{Reverse serialization: file to editable form}\label{sec:tf-reverse}

The reverse path is \code{load\_spec\_from\_file(path) -> dict}
(\file{pipeline/theory\_serialize.py:528}). Its job is to reconstruct the spec
dict from an existing \file{.theory.py} so the UI's ``Load'' button can
re-populate the form. The defining choice is that it reads the file's
\emph{abstract syntax tree} and never executes it.

\begin{defn}
The Python standard-library module \term{\code{ast}} (Abstract Syntax Tree)
parses source \emph{text} into a tree of node objects --- the program's grammar,
before execution. Inspecting this tree lets you learn what a program \emph{says}
without ever \emph{running} it, which is exactly what you want when reading a
configuration file you would rather not execute as a side effect of loading the
form.
\end{defn}

The whole reverse path is an \code{ast} walk. The load-bearing calls are
(\file{pipeline/theory\_serialize.py}, lines noted):

\begin{lstlisting}
import ast                                   # 545
tree = ast.parse(src)                        # 550  text -> AST
description = ast.get_docstring(tree)        # 553  module docstring
val = ast.literal_eval(node.value)           # 568  a *safe* eval of a literal
\end{lstlisting}

\code{ast.parse} turns the file into a module node (a malformed file raises
\code{SyntaxError} here). \code{ast.get\_docstring} pulls the leading
triple-quoted string. The critical one is \code{ast.literal\_eval}:

\begin{defn}
\term{\code{ast.literal\_eval}} evaluates a node \emph{only if it is a literal} ---
a number, string, list, dict, tuple, boolean, or \code{None}. It refuses function
calls and bare names, so it can \emph{never} run arbitrary code. This is why the
reverse path can safely recover a matrix default like \code{[[2,3],[1,3]]} or a
metadata tuple like \code{('dx', 1)} from the source text.
\end{defn}

The code wraps it in a local helper \code{\_lit(node, default=None)}
(\file{...:639}) that returns the fallback on
\code{ValueError}/\code{SyntaxError}/\code{TypeError}, and a sibling
\code{\_kwargs(call)} (\file{...:646}) that reads a call's keyword arguments into
a dict. With those in hand the walk proceeds:

\begin{enumerate}
  \item read the file and \code{ast.parse} it (\file{...:548--550});
  \item recover the description from the docstring, then strip the auto-header
        with \code{\_strip\_doc\_boilerplate} (\file{...:553--556});
  \item scan \code{tree.body} for the module-level \code{ast.Assign} nodes that
        target \code{DEFAULT\_FUNDAMENTAL} / \code{METADATA} and
        \code{literal\_eval} their values (\file{...:561--574});
  \item find the \code{build} function (an \code{ast.FunctionDef}), and its
        \code{ast.Return}; a missing \code{build()} or a return that is not a call
        raises \code{ValueError} (\file{...:577--585});
  \item \textbf{unwind the method chain}: starting at the return-call, step into
        \code{cur.func.value} while the node is an \code{ast.Call}, stopping at
        the constructor \code{ast.Name}; then \emph{reverse} so the constructor is
        first and methods follow in source order (\file{...:588--598});
  \item initialise the spec with all expected keys at their defaults, including
        the four toggle defaults that mirror the builder's own
        (\file{...:601--637});
  \item dispatch over the chain --- a long \code{if/elif} on the method name
        (\code{population}, \code{physical\_field}, \code{parameter},
        \code{define\_function}, \code{define\_kernel}, \code{declare\_cgf\_term},
        \code{set\_action\_text}, \code{set\_mf\_equation}, \code{equation}, the
        toggles, \code{boundary}/\code{initial}, \code{dyson}/\code{dyson\_order},
        \code{reference\_diffusion}, \code{build}) --- each branch reconstructing
        one spec sub-list (\file{...:649--865});
  \item sync \code{n\_populations} to \code{len(populations)} when the populations
        list is non-empty (\file{...:868--869}).
\end{enumerate}

A few dispatch branches deserve a sentence. The constructor branch accepts all
three class names (\file{...:651--652}). The \code{physical\_field} branch
\emph{inherits} a builder-level \code{\_default\_spatial\_dim} (set by an earlier
\code{.spatial\_dim(d)} call) when no per-field \code{spatial\_dim=} is given
(\file{...:682--683}) --- this is why the transient spec key from
\S\ref{sec:tf-spec} exists. The \code{declare\_cgf\_term} branch handles both the
legacy positional \code{response\_field} and the new \code{response\_legs} (which
it accepts as a list literal \emph{or} a comma-separated string). The toggle
branches use sensible no-argument defaults: \code{markovianize()} reads as
\code{True}, \code{operator\_ir()} reads as \code{True}.

Finally, \code{\_strip\_doc\_boilerplate(doc)} (\file{...:873}) removes the
auto-generated header so a reloaded theory shows only the user's own prose: it
drops the leading \code{"<name> --- text-driven theory file."} line (recognised by
its em-dash), stops at the ``Generated by\dots\ / Loaded by\dots'' lines, and trims
surrounding blanks. This is the same trimming \code{describe\_model} performed in
Chapter~\ref{ch:quickstart}.

\subsection{The round-trip guarantee and its limits}

The reverse path is an \emph{AST} round-trip, not a textual one, which fixes
exactly what survives:

\begin{itemize}
  \item String literals --- action text, equation right-hand sides, function and
        kernel expressions --- survive \emph{verbatim}.
  \item Python literals --- numeric defaults, \code{indexed\_by} lists, metadata
        tuples --- survive exactly via \code{literal\_eval}.
  \item Comments and exact formatting do \emph{not} survive.
\end{itemize}

\begin{gotcha}
\textbf{Non-literal arguments are lost on reload.} \code{\_lit} returns the
fallback for anything that is not \code{literal\_eval}-able. If a hand-written
theory passes a \emph{computed} expression as a builder keyword --- say
\code{default=2*0.4} instead of \code{default=0.8} --- then
\code{load\_spec\_from\_file} recovers \code{None} for it. Defaults are normally
plain literals, so this rarely bites, but a hand-authored file that puts logic in
its \code{build()} can carry meaning the AST round-trip silently flattens. The
loader (next section) \emph{executes} the file, so it sees the computed value; the
UI ``Load'' button, reading the AST, does not.
\end{gotcha}

\begin{gotcha}
\code{load\_spec\_from\_file} \emph{raises} on a malformed file: \code{ast.parse}
raises \code{SyntaxError}, and a missing \code{build()} or a \code{build()} that
does not return a call expression raises \code{ValueError}
(\file{...:581}, \file{...:585}). The UI must catch these.
\end{gotcha}

\section{Discovery and loading by name}

You already used the canonical loader in Chapter~\ref{ch:quickstart}. It lives in
\file{notebooks/daedalus.py} and is worth re-reading now that you know what it
loads. \code{daedalus.load\_theory(name)} (\file{notebooks/daedalus.py:68})
resolves \file{theories/<name>.theory.py}, imports it via \code{importlib.util}
--- \emph{executing} it --- and returns \code{(module.build(), module)}:

\begin{lstlisting}
spec = importlib.util.spec_from_file_location(f'theories.{name}', path)
mod = importlib.util.module_from_spec(spec)
spec.loader.exec_module(mod)
return mod.build(), mod
\end{lstlisting}

\code{spec\_from\_file\_location} builds an import ``spec'' for an arbitrary file
path (so the file need not be on \code{sys.path}); \code{module\_from\_spec}
creates an empty module object; \code{exec\_module} actually runs the file's
top-level code, defining \code{build}, \code{DEFAULT\_FUNDAMENTAL}, and
\code{METADATA}.

\begin{gotcha}
\textbf{The reverse path reads; the loader executes.} \code{load\_spec\_from\_file}
uses \code{ast} and never runs the file --- safe, but it cannot recover
non-literal values. \code{daedalus.load\_theory} (and the binary serializer's
\code{reload\_model}) \emph{do} execute the file. So a \file{.theory.py} runs real
Python at run time: a hand-authored theory may legitimately contain logic the
AST-only reverse path would flatten. If you put logic in \code{build()}, the
engine will honour it but the UI ``Load'' button may not faithfully reflect it.
\end{gotcha}

The companions are small. \code{list\_theories()} (\file{notebooks/daedalus.py:61})
lists every \file{theories/*.theory.py} and strips the suffix to give the loadable
name. \code{THEORIES\_DIR} (\file{...:55}) is \code{<repo\_root>/theories}, where
\code{repo\_root()} (\file{...:40}) walks up the directory tree until it finds a
\code{pipeline/} folder --- the same depth-robust walk the setup cell performed in
Chapter~\ref{ch:quickstart}.

\section{The binary cache (\texttt{saved\_theories/})}

We close where the chapter's opening note began: the \emph{other} thing called
serialization. When a theory is actually \emph{run}, the heavy symbolic objects
--- the propagator, the Taylor-expanded action, the per-diagram integrands --- are
pickled to disk under \file{saved\_theories/<slug>/} as SageMath \code{.sobj}
files plus a \code{manifest.json}. This cache is managed by
\file{pipeline/\_expand\_cache.py} and \file{msrjd/core/serialize.py} --- \emph{not}
by \code{theory\_serialize.py}.

\begin{defn}
\term{SageMath} is a large open-source computer-algebra system built on Python; it
provides exact symbolic objects (polynomials over chosen rings, symbolic matrices,
\dots). Its \code{save}/\code{load} functions pickle these objects to
\term{\code{.sobj}} files (the extension is appended automatically). You need Sage
to load a \code{.sobj}; you do \emph{not} need it to read a \file{.theory.py}.
\end{defn}

The on-disk \term{slug} is \code{model['name']} lowercased with non-alphanumerics
collapsed to underscores --- the \emph{same} pattern \code{\_slugify} uses for the
text-format filename (\file{pipeline/\_expand\_cache.py} uses an identical regex),
so the human-facing filename and the cache directory agree. Thus
\code{'OU Quartic (white noise)'} caches under
\file{saved\_theories/ou\_quartic\_white\_noise/}. The directory's contents:

\begin{center}
\begin{longtable}{@{}lll@{}}
\toprule
\textbf{File} & \textbf{Producer} & \textbf{Contents} \\
\midrule
\endhead
\code{manifest.json} & \code{\_expand\_cache.py} & index of cached artefacts
  (\code{key}, \code{stage}, \code{k}, \code{loop\_order}, \code{saved\_at}). \\
\code{propagator.sobj} & \code{\_expand\_cache.py} & Sage-pickled propagator
  (Taylor-order-independent). \\
\code{expand\_taylor<N>.sobj} & \code{\_expand\_cache.py} & Sage-pickled
  Taylor-expanded action at order $N$. \\
\code{unique\_typed\_mult\_v3\_...\_k<k>\_l<l>.sobj} &
  \code{\_expand\_cache.py} & Per-diagram integrand bundle, keyed by external legs,
  Taylor order, $k$, loop order. \\
\code{metadata.json} + \code{symbolic\_data.sobj} & \code{msrjd/core/serialize.py} &
  the older (separate) save format. \\
\bottomrule
\end{longtable}
\end{center}

A representative \code{manifest.json} entry (from
\file{saved\_theories/ou\_quartic\_white\_noise/}) shows how the cache key encodes
the request --- this is where the $k$ and $\ell$ from \code{METADATA}/\code{Config}
end up as file-name keys:

\begin{lstlisting}[style=console]
{ "key": "unique_typed_mult_v3_dx1_dx1_taylor6_k2_l2",
  "stage": "unique_typed_mult_v3_dx1_dx1_taylor6",
  "k": 2, "loop_order": 2,
  "saved_at": "2026-06-17T23:03:39.062450" }
\end{lstlisting}

\begin{gotcha}
\file{msrjd/core/serialize.py} defines its \emph{own} \code{save\_theory} /
\code{load\_theory} functions --- distinct from both \code{theory\_serialize.py}
and \code{daedalus.load\_theory}. They persist expanded symbolic objects, not
\file{.theory.py} text. Its \code{reload\_model} even uses raw \code{exec} to
re-import the model \code{.py}, so it runs arbitrary code. Three different
functions across the codebase are spelled ``\code{load\_theory}''; always check
which module you are in.
\end{gotcha}

\section{Authoring rules, in one place}

Most of the chapter's gotchas reduce to a short checklist for anyone hand-writing
or hand-editing a theory file:

\begin{enumerate}
  \item Write the action with \code{xt} for the response field and \code{\^{}} for
        powers; wrap per-population dynamics in \code{sum(... for i in pop)}.
  \item Declare the saddle with \code{.equation(lhs, rhs, population)} (preferred)
        or \code{.set\_mf\_equation('xstar', rhs)} (legacy). Never declare a
        \code{<natural>star} parameter --- the framework creates it.
  \item For a spatial model, give at least one physical field a non-zero
        \code{spatial\_dim} so the builder split picks
        \code{SpatialTheoryBuilder}; and if the action uses \code{Dt()/Lap()/Dx()}
        call syntax, \emph{keep} the \code{.operator\_ir()} line --- dropping it
        silently computes the wrong physics.
  \item Keep \code{recommended\_external\_fields} entries as \emph{tuples}.
  \item Prefer plain literals for builder keyword values; a computed expression
        works at run time but is lost when the UI reloads the file via the AST
        path.
  \item Remember the two serialization layers: a \file{.theory.py} is plain
        text you can read and diff; a \file{saved\_theories/<slug>/*.sobj} is a
        Sage pickle written when the model runs. Deleting the cache directory only
        forces a recompute; it never touches your source.
\end{enumerate}

With the model fully specified --- by hand or by form --- the next chapter opens up
the one feature this running example deliberately avoided: coloured noise and the
Markovian embedding that turns it back into a local-in-time field theory.
